% Part: first-order-logic
% Chapter: syntax-and-semantics
% Section: formation-sequences

\documentclass[../../../include/open-logic-section]{subfiles}

\begin{document}

\olfileid{pl}{syn}{fseq}

\olsection{جوړښتي لړۍ}

په استقرايي تعريف سره د !!{formula}s تعريفول، او د استقرا له لارې
د !!{formula}s د خاصيتونو د ثابتولو بشپړوونکې طريقه، ښکلې او اغېزمنه
لار ده۔ خو کله کله دا هم ګټوره وي چې د !!{formula}s جوړول له بېخه
پورته، ګام په ګام وڅېړو۔ دلته د \emph{جوړښتي لړۍ} په مفهوم همدا کار
کوو۔

\begin{defn}[د فارمولونو جوړښتي لړۍ]
\ollabel{defn:fseq-frm}
د ژبې~$\Lang L_0$ د نښو د توريزو لړيو يوه متناهي لړۍ
$\tuple{!A_0,\dotsc,!A_n}$ د $!A$ دپاره \emph{جوړښتي لړۍ} ده که
$!A \ident !A_n$ وي او د هر $i \leq n$ دپاره يا $!A_i$ اټومي
فارمول وي، يا داسې $j,k < i$ موجود وي چې له لاندې حالتونو څخه يو
صدق وکړي:
\begin{enumerate}
    \tagitem{prvNot}{$!A_i \ident \lnot !A_j$۔}{}%
    \tagitem{prvAnd}{$!A_i \ident (!A_j \land !A_k)$۔}{}%
    \tagitem{prvOr}{$!A_i \ident (!A_j \lor !A_k)$۔}{}%
    \tagitem{prvIf}{$!A_i \ident (!A_j \lif !A_k)$۔}{}%
    \tagitem{prvIff}{$!A_i \ident (!A_j \liff !A_k)$۔}{}%
\end{enumerate}
\end{defn}

\begin{ex}
\[
\tuple{
    \Obj p_0,
    \Obj p_1,
    (\Obj p_1 \land \Obj p_0),
    \lnot (\Obj p_1 \land \Obj p_0)
}
\]
د $\lnot (\Obj p_1 \land \Obj p_0)$ يوه جوړښتي لړۍ ده، او دا هم ده:
\[
\tuple{
    \Obj p_0,
    \Obj p_1,
    \Obj p_0,
    (\Obj p_1 \land \Obj p_0),
    (\Obj p_0 \lif \Obj p_1),
    \lnot (\Obj p_1 \land \Obj p_0)
}.
\]
%
لکه چې له دويمې بېلګې ښکاري، په جوړښتي لړيو کښې ``بې‌ګټې توکي''
هم راتلے شي: داسې فارمولونه چې زياتي وي يا په جوړونه کښې ونډه نۀ
لري۔
\end{ex}

\begin{prop}
\ollabel{prop:formed}
په~$\Frm[L_0]$ کښې هر !!{formula}~$!A$ يوه جوړښتي لړۍ لري۔
\end{prop}

\begin{proof}
فرض کړئ $!A$ اټومي دے۔ نو لړۍ~$\tuple{!A}$ د~$!A$ جوړښتي لړۍ ده۔
%
اوس فرض کړئ چې $!B$ او~$!C$ په ترتيب سره جوړښتي لړۍ
$\tuple{!B_0,\dotsc,!B_n}$ او $\tuple{!C_0,\dotsc,!C_m}$ لري۔
%
\begin{enumerate}
  \tagitem{prvNot}{که $!A \ident \lnot !B$ وي، نو
      $\tuple{!B_0,\dotsc,!B_n,\lnot !B_n}$
      د~$!A$ جوړښتي لړۍ ده۔}{}
  \tagitem{prvAnd}{که $!A \ident (!B \land !C)$ وي، نو
      $\tuple{!B_0,\dotsc,!B_n,!C_0,\dotsc,!C_m,(!B_n \land !C_m)}$
      د~$!A$ جوړښتي لړۍ ده۔}{}
  \tagitem{prvOr}{که $!A \ident (!B \lor !C)$ وي، نو
      $\tuple{!B_0,\dotsc,!B_n,!C_0,\dotsc,!C_m,(!B_n \lor !C_m)}$
      د~$!A$ جوړښتي لړۍ ده۔}{}
  \tagitem{prvIf}{که $!A \ident (!B \lif !C)$ وي، نو
      $\tuple{!B_0,\dotsc,!B_n,!C_0,\dotsc,!C_m,(!B_n \lif !C_m)}$
      د~$!A$ جوړښتي لړۍ ده۔}{}
  \tagitem{prvIff}{که $!A \ident (!B \liff !C)$ وي، نو
      $\tuple{!B_0,\dotsc,!B_n,!C_0,\dotsc,!C_m,(!B_n \liff !C_m)}$
      د~$!A$ جوړښتي لړۍ ده۔}{}
\end{enumerate}
د !!{formula}s دپاره د استقرا د اصل له مخې، هر !!{formula} يوه
جوړښتي لړۍ لري۔
\end{proof}

د دې عکس هم ثابتولے شو۔ دا ځکه مهم دے چې ښيي د فارمولونو د تعريف
زموږ دواړه لارې همارزښته دي: دواړه يو شان پايلې ورکوي۔ دا هم ترې
معلومېږي چې د جوړښتي لړيو پر اوږدوالي عادي استقرا کارولے شو او د
فارمولونو په اړه قضيې ثابتولے شو۔

\begin{lem}
\ollabel{lem:fseq-init}
فرض کړئ $\tuple{!A_0,\dotsc,!A_n}$ د~$!A_n$ جوړښتي لړۍ ده او
$k \leq n$ دے۔ نو $\tuple{!A_0,\dotsc,!A_k}$ د~$!A_k$ جوړښتي لړۍ
ده۔
\end{lem}

\begin{proof}
تمرين۔
\end{proof}

\begin{thm}
\ollabel{thm:fseq-frm-equiv}
$\Frm[L_0]$ د ژبې~$\Lang L_0$ د نښو د ټولو هغو توريزو لړيو سټ دے
چې يوه جوړښتي لړۍ لري۔
\end{thm}

\begin{proof}
$F$ دې د ژبې~$\Lang L_0$ د نښو د هغو ټولو توريزو لړيو سټ وي چې
جوړښتي لړۍ لري۔ په \olref[pl][syn][fseq]{prop:formed} کښې مو وليدل
چې $\Frm[L_0] \subseteq F$؛ اوس ئې عکس ثابتوو۔

فرض کړئ $!A$ يوه جوړښتي لړۍ $\tuple{!A_0,\dotsc,!A_n}$ لري۔ پر~$n$
د قوي استقرا له لارې ثابتوو چې $!A \in \Frm[L_0]$۔ زموږ استقرايي
فرض دا دے چې د نښو هره هغه توريزه لړۍ چې د $m < n$ اوږدوالي
جوړښتي لړۍ لري، په~$\Frm[L_0]$ کښې ده۔ د جوړښتي لړۍ د تعريف له مخې،
يا $!A_n$ اټومي دے، يا بايد داسې $j,k < n$ موجود وي چې له لاندې
حالتونو څخه يو صدق وکړي:
\begin{enumerate}
    \tagitem{prvNot}{$!A_n \ident \lnot !A_j$۔}{}%
    \tagitem{prvAnd}{$!A_n \ident (!A_j \land !A_k)$۔}{}%
    \tagitem{prvOr}{$!A_n \ident (!A_j \lor !A_k)$۔}{}%
    \tagitem{prvIf}{$!A_n \ident (!A_j \lif !A_k)$۔}{}%
    \tagitem{prvIff}{$!A_n \ident (!A_j \liff !A_k)$۔}{}%
\end{enumerate}
اوس حالت په حالت استدلال کوو۔ که $!A_n$ اټومي وي، نو
$!A_n \in \Frm[L_0]$۔ بل حالت کښې فرض کړئ چې $!A \ident
(!A_j \land !A_k)$۔ د \olref[pl][syn][fseq]{lem:fseq-init} له مخې،
$\tuple{!A_0,\dotsc,!A_j}$ او $\tuple{!A_0,\dotsc,!A_k}$ په ترتيب
سره د $!A_j$ او~$!A_k$ جوړښتي لړۍ دي۔ ځکه چې دا د~$!A$ د جوړښتي
لړۍ خاصې پيلنۍ فرعي لړۍ دي، د دواړو اوږدوالے له~$n$ څخه لږ دے۔ نو
د استقرايي فرض له مخې $!A_j$ او~$!A_k$ په~$\Frm[L_0]$ کښې دي؛ او
بيا د !!a{formula} د تعريف له مخې $(!A_j \land !A_k)$ هم پکې دے۔
نور حالتونه په ورته استدلال ثابتېږي۔
\textbf{د ژباړې د نسخې سمون (OLPL-002):} سرچينه په دې يوه کرښه کښې
د نحوي يوشانوالي د نښې پر ځاے د معنايي هم ارزښت نښه کاروي، حال دا
چې تعريف، ټول نور حالتونه او ثبوت د توريزو لړيو يوشانوالی غواړي۔
دلته د هماغه تعريف نښه بېرته کارول شوې ده۔
\end{proof}

\end{document}
